% !TEX root = main.tex
% ============================================================================
% CHAPTER 5: DISCUSSION AND CONCLUSION
% ============================================================================

\chapter{بحث و نتیجه‌گیری}
\section{مقدمه}

این فصل، بخش پایانی پژوهش حاضر را تشکیل می‌دهد و با نگاهی جامع به مسیر طی‌شده، از طرح مسئله تا ارزیابی تجربی، دستاوردها و یافته‌های اصلی را در قالبی منسجم ارائه می‌کند. ابتدا مروری اجمالی بر اهداف پژوهش صورت می‌گیرد، سپس میزان تحقق هر یک از اهداف بر پایه نتایج عددی فصل چهارم ارزیابی می‌شود. در ادامه، یافته‌های کلیدی، محدودیت‌های شناسایی‌شده و پیشنهادات مشخصی برای تحقیقات آتی ارائه می‌گردد.

پژوهش حاضر در پاسخ به یک نیاز واقعی و رو به رشد در حوزه امنیت اطلاعات تصویری شکل گرفت: طراحی الگوریتمی که نه تنها امنیت بالایی در برابر انواع حملات شناخته‌شده داشته باشد، بلکه از پایه نظری محکم مبتنی بر سیستم‌های آشوبی نمایی نسل جدید برخوردار باشد. الگوریتم پیشنهادی، با ترکیب سیستم سه‌بعدی \lr{Chen} برای جایگشت و ۹ سیستم آشوبی نمایی برای انتشار پویا، رویکردی چندلایه و نوآورانه به رمزنگاری تصاویر رنگی ارائه می‌دهد.

\section{مرور اهداف پژوهش}

در این بخش، هر یک از اهداف کلی و جزئی تعریف‌شده در فصل اول با نتایج حاصل از آزمایش‌های تجربی فصل چهارم مقایسه و میزان تحقق آن‌ها ارزیابی می‌شود.

\subsection{هدف کلی: طراحی و توسعه الگوریتم رمزنگاری}

هدف کلی پژوهش، طراحی و توسعه الگوریتمی مقاوم برای رمزنگاری تصاویر رنگی با بهره‌گیری از ترکیب ۹ سیستم آشوبی نمایی و سیستم پیوسته چن بود. این هدف به‌طور کامل محقق شده است: الگوریتمی پنج‌مرحله‌ای طراحی، پیاده‌سازی و ارزیابی شد که شامل جایگشت مبتنی بر \lr{Chen}، انتخاب پویای سیستم نمایی برای هر پیکسل، دو لایه \lr{XOR} مستقل و معکوس‌پذیری ۱۰۰\% می‌باشد. پیاده‌سازی کامل در محیط \pyver{} و اجرای موفق بر روی تصاویر استاندارد پایگاه‌های داده مرجع \lr{USC-SIPI} و \lr{Kodak}، تحقق این هدف را تأیید می‌کند.

\subsection{اهداف جزئی و میزان دستیابی}

همان‌گونه که در جدول ۵-۱ خلاصه شده است، اهداف جزئی پژوهش همراه با معیارهای ارزیابی و نتایج عددی متناظر با آن‌ها آورده شده است:

\begin{table}[htbp]
\centering
\small
\resizebox{\textwidth}{!}{%
\begin{tabular}{>{\raggedright\arraybackslash}p{0.32\textwidth}>{\raggedright\arraybackslash}p{0.34\textwidth}>{\raggedright\arraybackslash}p{0.34\textwidth}}
\hline
هدف جزئی & معیار ارزیابی & نتیجه حاصل \\
\hline
بررسی عملکرد ۹ سیستم نمایی در تولید دنباله تصادفی & آنتروپی شانون کانال‌های رمزشده & میانگین \lr{7.9973} از ۸ بیت ایده‌آل \\
بهبود مقاومت در برابر حملات آماری و جستجوی فراگیر & هیستوگرام یکنواخت، فضای کلید $> 2^{350}$ & هیستوگرام مسطح در تمام کانال‌ها؛ فضای کلید $\sim 10^{105}$ \\
افزایش پیچیدگی با انتخاب تصادفی سیستم برای هر پیکسل & \lr{NPCR} و ضریب همبستگی رمزشده & \lr{NPCR} تا \lr{99.48\%}؛ کاهش همبستگی از \lr{0.9989} به \lr{0.0082} \\
مقایسه با روش‌های مرجع در ادبیات & آنتروپی، \lr{NPCR}، \lr{UACI} در مقایسه با روش‌های مرجع & آنتروپی رقابت‌پذیر؛ \lr{NPCR} در محدوده مرجع؛ \lr{UACI} مطلوب در \lr{Kodak} \\
\hline
\end{tabular}%
}
\caption{ارزیابی میزان تحقق اهداف جزئی پژوهش بر اساس نتایج تجربی}
\label{tab:ch5_objectives}
\end{table}

\section{یافته‌های اصلی پژوهش}

بر اساس نتایج تجربی ارائه‌شده در فصل چهارم، یافته‌های اصلی این پژوهش را می‌توان در پنج محور اصلی دسته‌بندی کرد:

\subsection{معکوس‌پذیری کامل الگوریتم}

مهم‌ترین یافته عملیاتی پژوهش، تأیید معکوس‌پذیری ۱۰۰\% الگوریتم پیشنهادی است. مقدار $\mathrm{MSE}=0$ و $\mathrm{PSNR}=\infty$ برای تمام شش تصویر آزمایش‌شده نشان می‌دهد که هیچ‌گونه اطلاعاتی در فرآیند رمزنگاری و رمزگشایی از دست نمی‌رود. این ویژگی حیاتی، کاربرد الگوریتم را در سامانه‌های حساس نظیر ذخیره‌سازی تصاویر پزشکی یا انتقال داده‌های محرمانه ممکن می‌سازد.

\subsection{آنتروپی بالا و توزیع یکنواخت پیکسل‌ها}

میانگین آنتروپی شانون تصاویر رمزنگاری‌شده برابر با \lr{7.9973} بیت است که انحراف تنها \lr{0.0027} بیت از مقدار ایده‌آل ۸ بیت دارد. این نتیجه برای تمام تصاویر، صرف نظر از محتوا و ساختار آماری اولیه، به‌دست آمده است و نشان‌دهنده توانایی الگوریتم در یکسان‌سازی توزیع پیکسل‌ها مستقل از محتوای اولیه تصویر است.

\subsection{حذف مؤثر همبستگی مکانی}

ترکیب جایگشت مبتنی بر سیستم چن و انتشار \lr{XOR} دولایه، همبستگی مکانی پیکسل‌های مجاور را از میانگین \lr{0.9989} در تصاویر اصلی به \lr{0.0082} در تصاویر رمزنگاری‌شده کاهش داده است؛ یعنی کاهشی معادل ۹۹.۲\%. این نتیجه در هر سه جهت افقی، عمودی و قطری مشاهده شد.

\subsection{حساسیت بالا به کلید و مقاومت تفاضلی}

آزمون \lr{NPCR} با تغییر $10^{-15}$ در یکی از پارامترهای کلید نشان داد که تا بیش از ۹۹.۴۸\% از پیکسل‌های تصویر رمزنگاری‌شده تغییر می‌کنند. این رفتار تضمین می‌کند که حتی اگر مهاجم به کلیدی با خطای در حد دقت ماشین دسترسی داشته باشد، رمزگشایی صحیح ناممکن خواهد بود.

\subsection{فضای کلید بسیار بزرگ و مقاومت در برابر جستجوی فراگیر}

با هفت پارامتر پیوسته با دقت ۱۵ رقم اعشاری، فضای کلید الگوریتم برابر با $10^{105}$ است که معادل بیش از $2^{350}$ می‌باشد. این مقدار از آستانه امنیتی \lr{AES-256} (فضای کلید $2^{256}$) و آستانه \lr{NIST} ($2^{128}$) بسیار فراتر می‌رود و حمله جستجوی فراگیر را با هر سخت‌افزاری به لحاظ محاسباتی کاملاً ناممکن می‌سازد.

\section{ارزیابی فرضیه‌های پژوهش}

در این بخش، هر یک از فرضیه‌های مطرح‌شده در فصل اول با شواهد عددی حاصل از آزمایش‌ها بررسی می‌شود:

\begin{table}[htbp]
\centering
\small
\resizebox{\textwidth}{!}{%
\begin{tabular}{>{\raggedright\arraybackslash}p{0.44\textwidth}>{\raggedright\arraybackslash}p{0.38\textwidth}>{\raggedright\arraybackslash}p{0.18\textwidth}}
\hline
فرضیه پژوهش & شاهد عددی تجربی & وضعیت فرضیه \\
\hline
فرضیه ۱: ترکیب ۹ سیستم نمایی مقاومت الگوریتم را در برابر حملات آماری و تفاضلی افزایش می‌دهد & هیستوگرام نزدیک به یکنواخت در تمام تصاویر؛ آنتروپی میانگین = \lr{7.9973} بیت & تأیید شد \\
فرضیه ۲: جایگشت چن همبستگی پیکسل‌های مجاور را به صفر نزدیک می‌کند & کاهش همبستگی از \lr{0.9989} به \lr{0.0082} (کاهش ۹۹.۲\%) & تأیید شد \\
فرضیه ۳: انتخاب تصادفی نگاشت‌های نمایی، آنتروپی را به مقدار ایده‌آل نزدیک می‌کند & میانگین آنتروپی \lr{7.9973} بیت (انحراف تنها \lr{0.0027} از ایده‌آل) & تأیید شد \\
فرضیه ۴: ساختار پیشنهادی سرعت اجرای مناسبی برای کاربردهای عملی دارد & رمزنگاری \lr{22.7}\lr{---}\lr{25.7} ثانیه در پایتون تک‌هسته‌ای برای \imdim{512} & نیازمند بهینه‌سازی محاسباتی \\
فرضیه ۵: اعمال لایه انتشار دوم حساسیت به تغییرات جزئی کلید و تصویر را افزایش می‌دهد & \lr{NPCR} تا \lr{99.48\%} و هیستوگرام کاملاً مسطح & تأیید شد \\
\hline
\end{tabular}%
}
\caption{ارزیابی فرضیه‌های پژوهش بر اساس نتایج تجربی}
\label{tab:ch5_hypotheses}
\end{table}

\section{محدودیت‌های پژوهش}

محدودیت‌های اصلی شناسایی‌شده در این پژوهش عبارتند از:
\begin{itemize}
    \item زمان اجرای رمزنگاری در محیط مفسری پایتون به دلیل حلقه پردازش تک‌تک پیکسل‌ها مناسب کاربردهای بلادرنگ نیست؛ این محدودیت با برداری‌سازی یا پیاده‌سازی \lr{C++/CUDA} برطرف می‌شود.
    \item مقدار \lr{UACI} برای برخی تصاویر با محتوای یکدست (نظیر \lr{Airplane}) اندکی از مقدار ایده‌آل ۳۳.۴۶\% فاصله دارد که با افزودن مکانیزم بازخورد به مرحله انتشار قابل ارتقاء است.
    \item آزمون‌های آماری \lr{NIST SP 800-22} بر روی دنباله‌های تولیدشده در مراحل بعدی پژوهش باید تکمیل گردند.
\end{itemize}

\section{پیشنهادات برای تحقیقات آینده}

بر اساس یافته‌ها و حوزه‌های بهبود شناسایی‌شده، پیشنهادات زیر برای پژوهشگران آینده ارائه می‌شود:

\subsection{بهینه‌سازی کارایی محاسباتی و سرعت پردازش}
\begin{itemize}
    \item \textbf{برداری‌سازی کامل با \lr{NumPy}:} بازنویسی فرآیند انتخاب پویای نگاشت‌ها به صورت برداری برای دستیابی به شتاب ۱۰ تا ۵۰ برابری.
    \item \textbf{پیاده‌سازی موازی با \lr{CUDA} بر روی پردازنده گرافیکی (\lr{GPU}):} بهره‌گیری از ماهیت موازی‌سازی‌پذیر رمزنگاری پیکسل‌ها برای دستیابی به عملکرد بلادرنگ در ویدیوها و تصاویر فوق‌واضح.
    \item \textbf{توسعه ماژول‌های \lr{C/C++}:} کامپایل بخش‌های بحرانی حل معادلات با \lr{Cython} یا \lr{C++}.
\end{itemize}

\subsection{بهبود ساختار انتشار برای ارتقاء مقاومت تفاضلی (\lr{UACI})}
\begin{itemize}
    \item \textbf{افزودن مرحله بازخورد (\lr{Feedback}):} استفاده از مقدار رمزنگاری‌شده پیکسل قبلی به عنوان ورودی برای رمزنگاری پیکسل بعدی به منظور ایجاد وابستگی متوالی و افزایش یکپارچگی \lr{UACI}.
    \item \textbf{ترکیب عملگر \lr{XOR} با جمع ماژولار $\bmod 256$:} برای یکنواخت‌تر ساختن توزیع تغییرات شدت روشنایی در میان پیکسل‌ها.
\end{itemize}

\subsection{گسترش دامنه ارزیابی و کاربردهای نوین}
\begin{itemize}
    \item اجرای آزمون‌های رسمی تصادفی‌بودن \lr{NIST SP 800-22}.
    \item ارزیابی مقاومت الگوریتم بر روی مجموعه‌های داده مقیاس‌بزرگ نظیر \lr{ImageNet} و \lr{COCO} جهت تایید تعمیم‌پذیری آماری.
    \item تعمیم معماری به رمزنگاری تصاویر چندطیفی و ابرطیفی (\lr{Hyperspectral}).
    \item تلفیق الگوریتم با روش‌های پنهان‌نگاری و واترمارکینگ برای ایجاد سامانه‌های چندمنظوره اصالت‌سنجی و حفظ محرمانگی.
\end{itemize}

\section{نتیجه‌گیری نهایی}

این پژوهش نشان داد که بهره‌گیری از نگاشت‌های آشوبی نمایی نسل جدید با گستره آشوب پایدار وسیع، در ترکیب با سیستم‌های چندبعدی پیوسته، افقی نوین در طراحی الگوریتم‌های رمزنگاری تصویر چندلایه و مقاوم فراهم می‌سازد. الگوریتم پیشنهادی با موفقیت تمام مؤلفه‌های امنیتی شامل معکوس‌پذیری ۱۰۰\%، آنتروپی شانون بسیار بالا (\lr{7.9973} بیت)، حذف همبستگی مکانی، حساسیت تفاضلی بالا و فضای کلید ایمن $2^{350}$ را محقق ساخته و پاسخگوی نیازهای نوین امنیت داده‌های تصویری در فضای دیجیتال است.